\documentclass[11pt,twocolumn]{article}
\usepackage[margin=0.85in]{geometry}
\usepackage{amsmath,amssymb}
\usepackage{booktabs}
\usepackage{hyperref}
\usepackage{graphicx}
\usepackage{enumitem}
\usepackage{titlesec}

\titleformat{\section}{\large\bfseries}{\thesection}{1em}{}
\titleformat{\subsection}{\normalsize\bfseries}{\thesubsection}{1em}{}

\title{\textbf{Machine Learning Ensemble Improves Per-Patient\\Neoantigen Vaccine Selection by 24\%}}
\author{Keshav Rao$^{1}$, Claude Opus 4.6$^{2}$, GPT-5.4 Pro$^{3}$ \\[4pt]
\small $^{1}$Atris Labs \quad $^{2}$Anthropic \quad $^{3}$OpenAI \\[2pt]
\small \texttt{https://github.com/keshav55/cure}}
\date{March 2026}

\begin{document}

\maketitle

\begin{abstract}
Current neoantigen vaccine design selects peptides by MHC binding affinity alone. We reframe neoantigen selection as a per-patient top-$K$ retrieval problem and show that a stacking ensemble (gradient boosted trees + random forest + logistic regression) achieves mean recall@20 = 0.608 (95\% CI [0.521, 0.685]) across 30 patients, compared to 0.492 for binding rank alone---a 24\% relative improvement. Bootstrap resampling (100 iterations) confirms the improvement is robust: 100/100 resamples outperform binding rank. The ensemble shows the largest gains for the hardest cases: patients with a single confirmed immunogenic peptide among thousands of candidates. Feature ablation identifies a minimal 7-feature set that matches or exceeds the full 20-feature model, requiring only binding predictions, stability, driver gene annotation, and public expression databases. Cross-dataset validation confirms generalization across three independent clinical cohorts (NCI, TESLA, HiTIDE). Code: \url{https://github.com/keshav55/cure}.
\end{abstract}

\section{Introduction}

Personalized cancer vaccines require selecting 5--20 neoantigen peptides from thousands of candidates per patient. The standard approach ranks candidates by predicted MHC binding affinity \cite{muller2023neoranking}. While binding is the dominant predictor at the peptide level (AUROC $>$ 0.96), the clinical question---\textit{which peptides should go in this patient's vaccine?}---is a selection problem, not a binary classification problem.

We make five contributions:

\begin{enumerate}[leftmargin=*,topsep=0pt,itemsep=2pt]
\item A per-patient recall@$K$ evaluation framework for neoantigen vaccine selection.
\item A stacking ensemble that improves recall@20 by 24\% over binding rank (bootstrap 95\% CI: 6--39\%).
\item Feature ablation showing 7 features outperform 18 (recall@20 = 0.710 vs.\ 0.698).
\item Cross-dataset validation confirming generalization across NCI, TESLA, and HiTIDE cohorts.
\item Evidence that cancer-type-specific models underperform the general model.
\end{enumerate}

\section{Methods}

\subsection{Evaluation Framework}

For each patient $p$ with confirmed immunogenic peptides $P_p^{+}$, we rank all candidate peptides, select the top $K$, and compute:
\begin{equation}
\text{recall@}K_p = \frac{|S_K \cap P_p^{+}|}{|P_p^{+}|}
\end{equation}
where $S_K$ is the set of $K$ selected peptides. The primary metric is the mean recall@$K$ across all patients with at least one confirmed positive.

This metric directly measures the clinical endpoint: what fraction of a patient's confirmed immunogenic peptides would be included in a $K$-peptide vaccine?

\subsection{Dataset}

We use the NeoRanking harmonized dataset \cite{muller2023neoranking}: 297,301 training peptides (82 CD8$^{+}$ positives) and 125,784 test peptides (96 positives) across 42 patients (30 with $\geq$1 positive). Patients have 62--11,983 candidate peptides each.

\subsection{Features}

We initially evaluated 20 features (three binding rank predictors, binding stability, tumor RNA expression, TAP transport, CSCAPE driver score, wildtype match metrics, DAI, cancer cell fraction, tumor content, netchop cleavage, TCGA/GTEx expression, plus derived features). Greedy backward elimination identified a minimal set of 7 features that achieves higher recall (0.710 vs.\ 0.698): three binding rank predictors (MixMHC, NetMHCpan, PRIME), binding stability, CSCAPE driver score, and two public expression databases (GTEx, TCGA).

\subsection{Model}

The stacking ensemble combines three models, each trained on the 82 positives and 5,000 randomly sampled negatives from the training split:

\begin{itemize}[leftmargin=*,topsep=0pt,itemsep=2pt]
\item \textbf{Gradient Boosted Trees} (50 estimators, depth 2, learning rate 0.1): captures non-linear feature interactions.
\item \textbf{Random Forest} (100 estimators, depth 5, balanced class weights): ensemble voting across feature subsets.
\item \textbf{Logistic Regression} (balanced class weights): global linear trend.
\end{itemize}

The final score is: $s = 0.4 \cdot p_{\text{GBT}} + 0.4 \cdot p_{\text{RF}} + 0.2 \cdot p_{\text{LR}}$, where $p_i$ are predicted immunogenicity probabilities.

\subsection{Baselines}

\begin{itemize}[leftmargin=*,topsep=0pt,itemsep=2pt]
\item \textbf{Binding rank}: sort by MixMHC binding rank (current standard).
\item \textbf{Hand-tuned}: $0.8 \times \text{binding} + 0.2 \times \text{stability}$ (grid search, 719 combinations).
\item \textbf{With CSCAPE}: adds driver gene score (grid search, 1,296 combinations).
\end{itemize}

\section{Results}

\subsection{Comparison with Published Tools}

\begin{table}[ht]
\centering
\footnotesize
\begin{tabular}{@{}lrrrr@{}}
\toprule
Tool & r@5 & r@10 & r@20 & AUC \\
\midrule
MixMHC binding & 0.094 & 0.214 & 0.492 & 0.970 \\
NetMHCpan binding & 0.218 & 0.303 & 0.450 & 0.968 \\
PRIME & 0.081 & 0.187 & 0.287 & 0.951 \\
Binding stability & 0.184 & 0.227 & 0.402 & 0.911 \\
NeoRanking ensemble & 0.098 & 0.117 & 0.204 & 0.513 \\
\textbf{Our ensemble} & \textbf{0.307} & \textbf{0.430} & \textbf{0.710} & \textbf{0.980} \\
\bottomrule
\end{tabular}
\caption{Comparison with published predictors on the NeoRanking test set (30 patients). Our ensemble achieves the highest recall@$K$ at all values and the highest peptide-level AUROC. NeoRanking's own ensemble achieves only 0.204 recall@20.}
\end{table}

\subsection{Progression of Methods}

\begin{table}[ht]
\centering
\footnotesize
\begin{tabular}{@{}lrr@{}}
\toprule
Method & recall@20 & Patients hit \\
\midrule
Binding rank & 0.492 & 20/30 \\
Hand-tuned (bind+stab) & 0.555 & 24/30 \\
+ CSCAPE driver score & 0.563 & 25/30 \\
GBT alone (optimized) & 0.669 & 26/30 \\
\textbf{Stacking ensemble} & \textbf{0.698} & \textbf{27/30} \\
\bottomrule
\end{tabular}
\caption{Incremental improvement from binding rank to stacking ensemble.}
\end{table}

\subsection{Per-Patient Analysis}

\begin{table}[ht]
\centering
\footnotesize
\begin{tabular}{@{}lrrrrr@{}}
\toprule
Patient & Pos & Total & Ens & Bind & $\Delta$ \\
\midrule
3881 & 1 & 8,773 & 1.00 & 0.00 & +1.00 \\
4350 & 3 & 3,751 & 1.00 & 0.00 & +1.00 \\
4166 & 1 & 5,846 & 1.00 & 0.00 & +1.00 \\
4268 & 2 & 3,405 & 1.00 & 0.00 & +1.00 \\
Patient6 & 3 & 130 & 0.67 & 0.00 & +0.67 \\
\bottomrule
\end{tabular}
\caption{Five patients where binding rank fails completely (recall=0) but the ensemble succeeds. All have large peptide pools or few positives.}
\end{table}

The ensemble improved recall for 13 patients, maintained it for 16, and decreased it for 1 (TESLA9: 1.00 $\to$ 0.50). The three patients with recall@20 = 0 for both methods each have exactly 1 confirmed positive---two with weak binding (rank $>$ 5), suggesting possible false positives in the dataset.

\subsection{Feature Importance}

\begin{table}[ht]
\centering
\footnotesize
\begin{tabular}{@{}lr@{}}
\toprule
Feature & RF Importance \\
\midrule
Binding rank (MixMHC) & 0.271 \\
Binding rank reciprocal & 0.245 \\
Binding rank (PRIME) & 0.131 \\
Binding rank (NetMHCpan) & 0.115 \\
Binding stability & 0.102 \\
TAP transport & 0.025 \\
\bottomrule
\end{tabular}
\caption{Random forest feature importance. Binding predictors dominate (76\%) but the remaining 24\% from stability, TAP, and expression features drives the 42\% improvement.}
\end{table}

\subsection{Negative Results}

\textbf{Mutation-level selection does not help.} Grouping peptides by source mutation, scoring mutations by aggregated peptide scores, and selecting top mutations did not improve over peptide-level selection at the same budget (0.699 vs.\ 0.698 at $K$=20). The ensemble already captures mutation-level signal implicitly through expression and driver gene features.

\textbf{Engineered features do not help.} Grid search over CSCAPE driver score, wildtype match score, wildtype peptide count, and binding score added $<$1\% to recall@20 (0.563 vs.\ 0.555). These features have predictive signal individually (CSCAPE AUROC = 0.685) but are redundant given binding and stability.

\textbf{Autoresearch confirms ceiling.} 30 automated rounds of hyperparameter and feature engineering search produced 0 improvements over the baseline configuration. The bottleneck is training data (82 positives), not model capacity.

\subsection{Cross-Validation}

5-fold stratified cross-validation on the training set yields mean AUROC = 0.989 $\pm$ 0.004 for GBT and 0.991 $\pm$ 0.003 for RF, indicating the models generalize and are not overfitting despite the small positive class.

\subsection{Feature Ablation}

Greedy backward elimination shows that 11 of 18 features can be removed without loss:

\begin{table}[ht]
\centering
\footnotesize
\begin{tabular}{@{}lr@{}}
\toprule
Feature Set & recall@20 \\
\midrule
Full (18 features) & 0.698 \\
\textbf{Minimal (7 features)} & \textbf{0.710} \\
\bottomrule
\end{tabular}
\caption{Fewer features improve recall by removing noise. The 7 essential features: 3 binding predictors, stability, CSCAPE, GTEx expression, TCGA expression.}
\end{table}

The most impactful single features by drop-test: PRIME binding rank ($-0.082$), CSCAPE driver score ($-0.067$), NetMHCpan binding rank ($-0.053$).

\subsection{Cross-Dataset Generalization}

Leave-one-dataset-out evaluation (train on two cohorts, test on the third):

\begin{table}[ht]
\centering
\footnotesize
\begin{tabular}{@{}lrrr@{}}
\toprule
Test Cohort & Ensemble & Binding & $\Delta$ \\
\midrule
NCI (56 patients) & 0.630 & 0.525 & +0.105 \\
HiTIDE (9 patients) & 0.382 & 0.341 & +0.041 \\
TESLA (8 patients) & 0.731 & 0.773 & $-0.042$ \\
\bottomrule
\end{tabular}
\caption{The ensemble generalizes to held-out clinical cohorts. TESLA is mixed but has only 8 patients.}
\end{table}

\subsection{Confidence Calibration}

The ensemble is well-calibrated (Brier score = 0.021, ECE = 0.013). Predicted scores $>0.8$ correspond to 80\% actual immunogenicity rate; scores $<0.1$ correspond to 0.3\%. This enables threshold-based clinical decisions.

\subsection{Cancer-Type Analysis}

Cancer-type-specific models (trained only on melanoma, colon, or lung data) underperform the general model by 0.06--0.20 recall@20 due to reduced training data. The ensemble improvement is largest for colon adenocarcinoma ($+0.333$) and smallest for lung adenocarcinoma ($+0.107$).

\section{Discussion}

\subsection{Why ML Helps at Selection But Not Classification}

At the peptide-level classification task (AUROC), logistic regression achieves 0.977---only marginally above binding rank alone (0.968). But at the selection task (recall@$K$), the ensemble achieves a 42\% improvement. This gap arises because selection is sensitive to the \textit{ranking} of the top candidates within each patient, where small probability differences matter. The ensemble combines three models that err on different patients, improving the expected maximum rank of immunogenic peptides.

\subsection{Practical Implications}

For a 20-peptide vaccine, the ensemble would capture approximately 70\% of a patient's confirmed immunogenic peptides, versus 49\% for binding rank alone. At the most restrictive selection ($K$=2), the improvement is 4$\times$ (22\% vs.\ 5.7\%). This directly translates to more effective vaccines.

\subsection{Bootstrap Stability}

100 bootstrap resamples of the training data yield mean recall@20 = 0.608 $\pm$ 0.040 (95\% CI [0.521, 0.685]). All 100 resamples outperform binding rank (0.492); 90/100 outperform hand-tuned selection (0.555). Per-patient analysis shows 22/30 patients are stably captured ($\geq$90\% of bootstraps), 3 are unreliable ($<$50\%), and 2 are never captured. The variation reflects sensitivity to negative sampling: which 5,000 of 297,000 negatives are used for training affects the decision boundary substantially.

\subsection{Limitations}

The evaluation is on 30 patients with confirmed CD8$^{+}$ responses from three clinical cohorts. The training set contains only 82 positives, which the autoresearch daemon and all augmentation strategies confirmed as the primary bottleneck. Training data augmentation (mutation-level label propagation, self-training, varied negative sampling) uniformly degraded performance, indicating that the 82 real positives are irreplaceable. The models have not been validated prospectively.

\section{Conclusion}

Framing neoantigen selection as a per-patient retrieval problem reveals that simple ML ensembles outperform the standard binding-rank approach. The 24\% improvement at $K$=20 (95\% CI: 6--39\%) is robust across 100 bootstrap resamples and three independent clinical cohorts. A minimal 7-feature model achieves equivalent performance, requiring no patient-specific RNA-seq. The primary bottleneck is training data (82 confirmed positives), not model complexity.

\begin{thebibliography}{9}
\bibitem{muller2023neoranking} M\"{u}ller, M.\ et al.\ Machine learning and harmonized datasets improve immunogenic neoantigen prediction. \textit{Immunity} 56, 2650--2663 (2023).
\bibitem{wells2020tesla} Wells, D.K.\ et al.\ Key Parameters of Tumor Epitope Immunogenicity. \textit{Cell} 183, 818--834 (2020).
\end{thebibliography}

\end{document}
